% Source: 52_Chapter_40_A_Bridge_for_Intelligent_Moderns_Protes.txt
\bookchapter[A Bridge for Intelligent Moderns]{CHAPTER 40}{A BRIDGE FOR INTELLIGENT MODERNS}{Protestant Liberalism}

\leadin{O}{N HIS EIGHTY-FIFTH BIRTHDAY} in 1920, Lyman Abbott, who had been one of America’s most
influential ministers in the 1890s, looked back three-quarters of a century to his staunch Puritan
upbringing. He recalled his youthful view of God as “a kind of awful and omnipresent police justice”
and his own self-image as “a scared culprit who knows he is liable to punishment but does not
clearly know why.”

Long before 1920, however, along with many other Americans and Europeans, Abbott had ceased to think
of God as an “omnipresent policeman” and man as a “scared culprit.” The Western world had undergone
too many changes and adopted too many new ideas in the last quarter of the nineteenth century.

Abbott was typical of a large number of American ministers whose background included a pious
Protestant home but who had studied in Germany or in an American seminary where continental
scholarship was treasured and who had adopted “liberal” religious convictions.

The events of the twentieth century have been unkind to the liberal creed, but every major
Protestant denomination continues to reflect the impact of liberal theology. It is hard to argue
with Professor Sydney E. Ahlstrom’s judgment when he says the liberals “precipitated the most
fundamental controversy to wrack the churches since the age of the Reformation.” The reason lies in
their ambitious objective. They tried to lead the Protestant churches into the new world of modern
science, modern philosophy, and modern history. In his autobiography The Living of These Days, Harry
Emerson Fosdick, minister at the influential Riverside Church in New York City, put it well when he
said the central aim of liberal theology was to make it possible for a man “to be both an
intelligent modern and a serious Christian.”

\section{THE AIMS OF PROTESTANT LIBERALISM}

Protestant liberalism, then, engaged a problem as old as Christianity itself: how do Christians make
their faith meaningful in a new world of thought without distorting or destroying the gospel? The
apostle Paul tried and succeeded. The early Gnostics tried and failed. The jury is still out on
liberalism, but Christian public opinion tilts heavily in the direction of failure. No one expressed
the irony of liberalism better than H. Richard Niebuhr when he said in liberalism “a God without
wrath brought men without sin into a kingdom without judgment through the ministrations of a Christ
without a Cross.”

Beyond its rather clear goal, definitions of religious liberalism are as varied as those of
political liberalism. Many deny that Protestant liberalism is a theology. They prefer an outlook or
an approach or a spirit. Thus Henry Sloane Coffin at New York’s Union Seminary once said liberalism
is that spirit that reveres truth supremely and therefore craves freedom to discuss, to publish, and
to pursue what it believes to be true.

\begin{sidebar}{Theological Liberalism}
Theological Liberalism promotes a progressive ethical vision derived from the Christian message,
while denying Christianity’s traditional understanding of doctrine.

Liberalism’s perfect exemplar is Adolf von Harnack (1851–1930). He made pioneering contributions to
New Testament studies, theology, ethics, and historical theology (a field he revolutionized). He
argued that Jesus had been misunderstood and misrepresented in most Christian sources, including the
Gospels. Jesus had been a simple Hebrew prophet calling people to moral awakening. The message of
the genuine Jesus could be recovered, however, once the corrupt layers of tradition had been
removed, extracting husk from the kernel of enduring value.

According to Harnack, Jesus taught that the kingdom of God was not a coming supernatural cataclysm
but an ethical kingdom. He taught about the fatherhood of God and the brotherhood of mankind; each
and every soul was of infinite value. Jesus also called his followers to a higher righteousness
(beyond that of the Pharisees) that was rooted in God’s command to love.
\end{sidebar}

No doubt this is the outlook of liberals, but is that all? Doesn’t this spirit lead to identifiable
convictions? I think so. And that spirit and those convictions together constitute Protestant
liberalism.

It might be helpful to think of liberal theology as a suspension bridge. The footing of one tower is
planted upon modern thought and the foundation of the other rests upon Christian experience.
Unfortunately, the ground around both towers is shifting soil, and those who take the bridge
disagree over which is the safer side. That is why Professor Kenneth Cauthen finds two fundamental
types of liberalism. He calls them “evangelical liberalism” and “modernistic liberalism.”

Cauthen suggests that the evangelical liberals were “serious Christians,” to borrow Fosdick’s terms,
who were searching for a theology that could serve “intelligent moderns.” Evangelical liberals,
then, took greater confidence in the tower resting on Christian experience.

On the other side of the bridge were modernistic liberals who were intelligent moderns hoping to be
considered serious Christians in some sense. They found greater support in the tower resting on
modern thought.

Perhaps the best way to explore theological liberalism, then, is to take a close look at modern
thought and then at Christian experience.

Liberals believed that Christian theology had to come to terms with modern science if it ever hoped
to claim and hold the allegiance of intelligent men of the day. They refused, therefore, to accept
religious beliefs on authority alone. They insisted that faith had to pass the tests of reason and
experience. Man’s mind, they believed, was capable of thinking God’s thoughts after him, and the
best clues to the nature of God were human intuition and reason.

The Christian, they said, should keep his mind open to truth from any source. New facts may well
change traditional beliefs that rest on no more than custom and time, but unexamined faith is not
worth having.

By surrendering so completely to the modern mind liberals accepted the assumption that the universe
was one grand, harmonious machine, or perhaps an extremely complex growing organism. Whatever the
image, a watch or a plant, the point is unity, harmony, coherence.

The biblical account of creation, however, recognizes certain important orders in the universe:
inanimate matter, plants, animals, man, and God. That didn’t bother liberal theology. It pressed on
for unity or continuity. It reduced distinctions between revelation and natural religion, between
Christianity and other religions, between saved and lost, between Christ and other men, between man
and God.

Two technical theological terms are crucial here: immanence and transcendence. Immanence carries the
idea of God’s presence in the world and working through nature. An extreme version of immanence is
found in pantheism, which claims that God and the world are identical. Transcendence implies the
reality of God apart from the world. An extreme version of transcendence is found among the deists,
for whom God is as separate from the world as a watchmaker from his watch.

Liberals felt that the old orthodox Christian idea of a God somewhere beyond the universe was
unacceptable to modern men. Some liberals criticized seeing God as a distant and remote being,
though this view may have more to do with deism than genuine Christian theology. Opposed to a remote
understanding of God, they linked the divine with mankind’s conscious being. The life coursing
through nature and man they called God. This immanent view of God seemed to fit the results of
scientific studies. Instead of suddenly breaking through the clouds to create the world, God, they
said, had been working for ages through natural law, slowly building the universe as we find it
today. Most liberals agreed with the poet who said, “Some call it evolution, and others call it
God.”

Evolution was the theory that held that all complex living things have developed from simple forms
through the operation of natural selection. Thus no species is fixed and changeless. In 1785 James
Hutton had attributed the earth’s development to natural rather than supernatural causes.
Confirmation of the view came in Sir Charles Lyell’s epoch-making Principles of Geology (1830).
Lyell showed that the earth’s surface had been formed by natural causes operating over a vast period
of time. Such a conception of geologic time was essential for any theory of evolution based on
changes in species over many thousands of generations.

\section{ENTER CHARLES DARWIN}

The scientist whose name became synonymous with evolution was Charles Darwin (1809–1882). After
studying medicine and preparing at Cambridge University for the ministry, Darwin became a
naturalist. From 1831 to 1836 he studied the specimens he had collected while on a surveying
expedition with the ship Beagle along the coast of South America.

In 1859 Darwin’s views appeared in his Origin of Species. He contended “that species have been
modified during a long course of descent~.~.~. chiefly through the normal selection of numerous
successive, slight, favourable variations.” The Origin of Species, the most important book of the
century, revolutionized the concepts about the origin and evolution of life on planet earth. Darwin
followed his first bombshell by a second. In 1871 his Descent of Man applied the natural selection
to human beings and reached the controversial conclusion that man’s ancestors were probably
monkey-like animals.

Such conclusions threw many religious people on the defensive. Some vigorously rejected the new
scientific views. If man is not specially created by God and fallen from God’s favor, where is the
need for Christ’s salvation? Others attempted to reconcile their religious beliefs with evolution.
As time went on liberals came to believe that the evolutionary theory supplemented rather than
contradicted the basics of Christianity. They considered growth and development as God’s way of
revealing himself to man. In 1892 Lyman Abbott, then minister at the Plymouth Church in Brooklyn,
New York, wrote The Evolution of Christianity and attempted to show that “in the spiritual, as in
the physical, God is the secret and source of light.” He spoke of the evolution of the Bible, of the
church, and even of the soul.

As serious as the challenge of science was to orthodox Christianity, it was clearly secondary to the
new views of history. Science could only question God’s rule in the physical world, but historical
criticism advanced directly to the domain of the Christian faith, to the revelation of God in the
Bible.

The term for the application of the principles of history to the Bible is biblical criticism. The
term criticism may be somewhat misleading. It is true that many critical scholars reach conclusions
that challenge the Bible’s veracity and value as a witness to God’s purpose. But the term actually
identifies a strategy for knowing and learning. A scholar is critical if he or she subjects the
material to an inspection of careful (critical) scrutiny. For example, a modern or critical scholar
may inspect the New Testament materials to see if there actually is historical evidence to support
the church’s teaching that the Son was God or if this doctrine rests upon tradition. Although less
believable in our day, this movement assumed that it was immune from corrupting traditions itself.

Biblical criticism came to be expressed in two forms, what is sometimes called lower and higher
criticism. Lower criticism tried to reconstruct the precise wording of the originals (presently
lost) from an inspection of the many surviving ancient manuscripts. This scholarly pursuit may alarm
conservative Christians initially; but conservatives have contributed significantly to this
endeavor. After all, they have a great interest in the discernment and discovery of the exact
wording of the originals.

Higher criticism, however, proved another matter. The higher critic is not primarily interested in
the original wording of the text; he is interested in the meaning of the text. The higher critic
wants to inspect the text historically to see if the event which it reports really happened. To
achieve this, he studies the background, authorship, and setting of the text. The higher critic
believes that we can only understand the Bible if we see it against its background. The methods of
higher criticism were not new, but they had been limited to writings other than Scripture. In the
nineteenth century the method was applied to the Bible as if the Bible were any other ancient book
whose credentials had to pass the standards of historical methods. In actual practice the critics
did not apply critical methodology as they would to any other book. Many held assumptions that God
could not act in history in miraculous or routine ways. Others assumed the Bible story and teaching
to be untrue unless it could be shown true by an independent source beyond the Bible. This did
result in conclusions that shook orthodoxy.

Several critical conclusions troubled believers. Critics believed that the books of Moses were
actually the product of a centuries-long composition that included four great literary traditions
being combined throughout Israel’s history. Some critics speculated that some prophecies seemed true
because they were actually written after the events occurred. Others challenged the accuracy of the
Gospels, especially John. One of the central concerns of higher criticism was the search for the
“historical Jesus.” The critics assumed that Jesus, as he lived in history, was different from the
Jesus whom we find portrayed in the Gospels. They assumed that the early church and the gospel
writers had added many things to the biblical account so that the problem was to sift the authentic
sayings and doings of Jesus from the later additions.

The work of Strauss (1809–1874) rejected the previous approaches. He held that it was no longer
possible to embrace a straightforward supernaturalist approach or the face-saving approach of
rationalism. The rationalists held that readers could discover rational explanations for Jesus’
miracles that the gospel writers misunderstood; for example, Jesus was walking on a jetty and not on
the water. Also unsatisfactory was the notion that the gospel writers were intentionally fraudulent
in creating Jesus’ story. Strauss argued that the gospel’s representation of Jesus reflects a
mythological perspective of the ancients.

Myths were the vehicles that ancient people employed to grasp great people and events; the awareness
of the myth-shaped understanding of the world gave the modern reader the necessary insight to
appropriately assess the important ideas carried in the ancient myth stories.

\section{THE IMPACT OF BIBLICAL CRITICISM}

As important, however, than any of the particular conclusions of biblical critics was the fact that
criticism threw into question the general trustworthiness of the Bible. Liberals welcomed higher
criticism because they recognized a radically different view of the Bible was necessary for
intelligent moderns. They felt free from the need to apologize for the whole Bible as the infallible
Word of God. They no longer had to defend a God who ordered the Israelites to kill their enemies to
the last woman and child or who sent bears to maul children who poked fun at a prophet.

The studies of the higher critics, said the liberals, make it clear that God has revealed himself
through an evolutionary process. It began with primitive, bloodthirsty ideas and showed how the Jews
slowly came to grasp the idea of a righteous God who can be served only by one who does justly,
loves mercy, and walks humbly with his God. This evolutionary revelation of God finds its
fulfillment in Jesus, where God is portrayed as the loving Father of all men.

When liberals could no longer rest in the traditional doctrines of orthodoxy, which they felt
science and history had destroyed, they found their needed assurance in the other pillar of their
bridge: Christian experience.

In the early nineteenth century an artistic and intellectual movement arose called Romanticism.
While it is popularly understood as a protest against the overly rationalistic mindset of the
Enlightenment, Romanticism discerned a dynamic energy and force surging within nature. This strong
vitality was not exactly irrational but rather beyond reason. Romantic paintings often represent
this vital force as a great storm that is at once dangerous and fascinating to the onlooker. People
could discern this and be swept up into its creative and mental powers. Encountering and engaging
this vitality came to be of first importance. It was engaged by an intuitive awareness or sense and
not as a logical deduction. Why trouble about formal and external creeds when so intimate and so
undeniable a certainty ruled every soul? As Tennyson wrote:

\begin{poem}
Speak to Him, thou, for He hears, and Spirit can meet—\\
Closer is He than breathing, and nearer than hands and feet.
\end{poem}

When he was asked to interpret the American situation before the International Congregational
Council in London in 1891, Dr. Lewis F. Stearns of Bangor Seminary in Maine said,

\begin{blockquote}
We are coming to understand that it is the recognition of the invincible reality of spiritual
Christianity which is going to give our theology its great power in the future~.~.~.~. Criticism may
assail the historical facts of revelation: rationalism may urge objections to its doctrines; but the
surf on our coast of Maine might as easily overthrow the granite cliffs against which it breaks as
criticism and rationalism disturb the Christian realities which stand firm in the experience of the
individual believer and the church.
\end{blockquote}

\section{INFLUENTIAL SPOKESMEN FOR LIBERALISM}

The two most influential spokesmen for Christian experience were Friedrich Schleiermacher
(1768–1834) and Albrecht Ritschl (1822–1889). Both were German theologians. Schleiermacher taught at
the newly founded University of Berlin and Ritschl at Bonn and Gottingen.

According to Professor William Hordern, Schleiermacher attempted to rehabilitate religion among
intellectuals, insisting that the great debates over proofs of God and the abstract doctrinal
descriptions of faith were, at best, a secondary expression of religion. At the heart of religion
was an awareness of our absolute dependence and vulnerability before the grandeur of God. He
appealed to sophisticated intellectuals that they were right to reject the abstract versions of
Christianity seen in Protestant Scholasticism. In fact, their romantic sensibilities were inherently
religious, and Christianity, rightly understood, accounts for this experience better than any other
religion. After Schleiermacher, the idea became commonplace that religions are differing expressions
of the same inner religious experience shared by all people. The uniqueness of Christ, says
Schleiermacher, is not in some doctrine about Jesus or in some miraculous origin such as the Virgin
Birth: “The real miracle is Jesus himself. In Jesus we find a man who had the sense of
God-consciousness to a supreme degree.” Where we all have flashes of God, he had complete knowledge.
Where we give fitful obedience, he gave complete obedience. “As the Godfilled man,” Jesus was “our
great pioneer” in the realm of the spirit and morality.

Since Jesus has the full and complete knowledge of God, he is able to communicate the consciousness
of God to others. Through Jesus we can come into a vital and living relationship with God. The
church is the living witness to the fact that down through the centuries men have come to a vital
God consciousness through their contact with the life of Jesus. This leads to a true reunion with
our fellow men in brotherly living.

Schleiermacher was “the father of modern theology” primarily because he shifted the basis of the
Christian faith from the Bible to “religious experience.” Albrecht Ritschl, the most influential
theologian in the late nineteenth century and the principal teacher of American liberals, focused
religious experience upon the historical Jesus.

To Ritschl religion had to be practical. It must begin with the question, What must I do to be
saved? But if that question means, How can I go to heaven when I die? then it is a theoretical
question. To be saved means to live a new life, to be saved from sin, selfishness, fear, and guilt.

According to Hordern, “To be practical, Christianity must be built upon fact, so Ritschl welcomed
the search for the historical Jesus. The great Christian fact is the impact that Jesus has made upon
the church through the centuries.” Nature cannot introduce us to God because it speaks ambiguously
about God. “We find God instead in history, where movements arise dedicated to the values that make
life meaningful. The task of theology is to turn men again to Jesus and remind them anew of what it
means to follow him.”

For Ritschl religion rests upon the values of men, not upon the truth of science. Science tells us
the facts, things as they are; but religion weighs the facts and counts some more valuable than
others. “The great fact about man is that, although he is a product of nature and evolution, he has
a sense of values.” We can explain this only if we recognize that the universe creates not only
atoms and molecules but also values. “God is the necessary postulate to explain this sense of worth
in man.”

Many Christians in the late nineteenth century, says Hordern, found Ritschl’s approach helpful. It
appeared to free the Christian faith from the destructive impact of history and science. It allowed
biblical criticism to take the way of science: to decide the facts about authorship, date, and
meaning of the biblical books. But it also recognized that religion is much more than facts. Science
cannot rule on values, the stuff of religion.

If biblical criticism denies Jesus’ miracles, his Virgin Birth, his preexistence, this does not make
Jesus less valuable to us. Belief in the divinity of Jesus does not rest in any of these; it rests
solely on the fact that he is the source of a value-creating movement; he has led men to find the
God of values. That is, Jesus’ life was the embodiment of such high ethical ideals and attainment
that we are inspired to live as he did. Jesus is divine in the sense that he can do for us what God
does; he makes us conscious of the highest in life. From Jesus’ influence, then, comes the church, a
value-creating community: the spearhead for building a society inspired by love, the kingdom of God
upon earth.

The impact of liberalism was not limited to any single denomination or country. It challenged
traditional orthodox bodies all over Europe and North America. We may take the Congregationalists in
the United States as an example of many other Christian groups.

\section{AN EXAMPLE OF LIBERAL CHANGE}

Liberal theology appeared among New England churches under the title New Theology. The leading
advocates of the New Theology emerged from within traditional New England Calvinism. Theodore
Thornton Munger was minister of the United Church (Congregational) in New Haven, Connecticut. Newman
Smyth, his colleague in New Haven, served the Center Church (Congregational) for twenty-five years.
George Angier Gordon was minister of historic Old South Church in Boston. Washington Gladden early
served in New England but found his greatest success at the First Congregational Church in Columbus,
Ohio. George Harris was Abbott Professor of Theology at Andover Theological Seminary and, supported
by his colleagues, was spokesman for the New Theology through The Andover Review. Finally, there was
the great popularizer of the movement, Lyman Abbott. “Probably no man,” said F. H. Foster, “ever
wielded in America so powerful an influence in the direction of a liberal theology.”

Scores of men clustered around these early leaders—Henry Ward Beecher, Egbert Smyth, William Jewett
Tucker, Lewis French Stearns, William Newton Clarke, and others—but the character of the movement
remained basically the same. It was a protest in the interests of modern thought against the “old
theology” of evangelical Puritanism, usually called New England Theology.

In 1881 Edwards Amasa Park, perhaps the last outstanding spokesman for New England Theology,
resigned from the influential Abbot Professorship of Theology at Andover. Two years later Harris
succeeded him. We may consider that event the watershed of the old and the new.

Prior to 1880 most New England ministers held to the sovereignty of God; to the innate depravity of
mankind (which they traced to the sin of the first man); to the atonement of Jesus Christ, the
ground of man’s forgiveness of sin; to the Holy Spirit as essential to conversion; and to the
eternal separation of the saved and lost in heaven or hell.

After 1880 every one of these beliefs came under heated fire from the liberals. The most celebrated
controversy erupted at Andover Seminary. The seminary had been established by New England
Congregationalists in 1808 to counter the Unitarian tendencies of Harvard. Attempting to preserve
Andover’s orthodoxy, the founders required faculty subscription to a creed summarizing their
inherited Calvinism. By 1880, however, faculty members, under the influence of liberalism, found
this requirement impossible and said so. The spark that lit the flames of controversy was a series
of articles in the Andover Review. Egbert C. Smyth, William Jewett Tucker, and George Harris of the
faculty argued that heathen who die without any knowledge of the gospel will have an opportunity in
the future life either to accept or to reject the gospel before facing final judgment. Step by step
in the debate that followed faculty members moved toward the public defense of liberal theology.

The Board of Visitors, one of the ruling bodies of the seminary, finally brought charges against
Smyth for his departure from the Creed: a kind of test case. After years of moves and counter-moves,
the Supreme Court of Massachusetts in 1892 voided the action of the Board of Visitors in dismissing
Smyth. By that time almost every American denomination had on its hands its own celebrated case of
heresy.

\readinghead{Suggestions for Further Reading}
\begin{readinglist}
\entry Ahlstrom, Sydney E. A Religious History of the American People. New Haven, CT: Yale
  University Press, 1972.
\entry Cauthen, Kenneth. The Impact of American Religious Liberalism. New York: Harper \& Row, 1962.
\entry Dorrien, Gary J. The Making of American Liberal Theology, 1900–1950. Louisville, KY:
  Westminster John Knox Press, 2003.
\entry Hordern, William E. A Layman’s Guide to Protestant Theology. New York: MacMillan, 1968.
\entry Shriver, George H. American Religious Heretics. Nashville: Abingdon Press, 1966.
\end{readinglist}
